\chapter{翻前 / Pre-flop}

\section{本章导入}

本章将围绕翻前决策展开。翻前是整手牌的起点，位置、起手牌范围和主动权会直接影响后续每条街的决策质量。初学者常见的问题不是翻后完全不会打，而是翻前已经用过宽、过弱、过被动的范围进入了困难局面。

翻前学习的目标，是建立清晰的入池纪律：什么位置可以 open raise，哪些牌适合 call，什么时候考虑 3-bet，什么时候应当直接 fold。

\section{核心概念}

本章将介绍位置、起手牌范围（range）、open raise、call、3-bet、fold 和主动权。读者需要理解，翻前不是孤立地评价两张手牌，而是结合位置、对手类型和后续可操作性进行判断。

\section{新手常见误区}

娱乐玩家常见错误包括入池范围过宽、喜欢 limp、面对加注过度跟注、低估位置劣势，以及拿边缘牌制造复杂局面。本节后续将逐项说明这些错误如何放大翻后压力。

\section{实战思考框架}

翻前可以按照“位置--手牌--前序行动--对手倾向--后续计划”的顺序思考。每一次入池都应当有理由，而不是因为“这手牌看起来还可以”。

\section{牌例拆解}

\subsection{牌例一：待补充}

本牌例将用于分析早期位置用边缘牌入池带来的连锁问题。

\subsection{牌例二：待补充}

本牌例将用于比较 call 与 3-bet 在不同对手类型下的策略差异。

\section{本章总结}

翻前质量决定了整手牌的起点。新手应先建立紧凑、清晰、有位置意识的范围，再逐步学习更复杂的调整。

\section{课后训练}

请为 UTG、CO、BTN 三个位置分别列出你愿意率先加注的起手牌范围，并标记其中最容易误玩的边缘牌。
